\section{模型成为新一代 OS}

访谈提出“模型是新一代操作系统”。这不是说模型替代 Linux、Windows 或 iOS，而是说模型可能成为用户意图和应用执行之间的新调度层。过去 OS 管理硬件、文件、进程和应用；未来模型/Agent OS 可能管理意图、工具、上下文、权限和任务状态。

\begin{figure}[H]
\centering
\includegraphics[width=0.92\textwidth]{model-as-os.png}
\caption{模型作为新一代 OS：意图理解、工具调度和状态记忆。}
\end{figure}

\begin{knowledgebox}{读图：Model OS 管什么}
图中模型 OS 位于中心，连接用户意图、应用服务、数据记忆、执行环境和安全权限。它支持的判断是：模型不只是一个 app，而可能成为新应用入口和任务调度层。
\end{knowledgebox}

\subsection{为什么 Coding 是 OS 化入口}

代码是数字世界最强的形式语言。模型如果能写代码、改代码、调用 API、运行脚本，就能跨越单个 app 的边界，成为操作多个系统的中间层。Coding 因此是模型 OS 化的关键入口。

\begin{importantbox}{从 app 到 OS 的变化}
App 是单一功能入口；OS 是资源和任务调度层。模型若能理解意图、调用工具、管理上下文和执行任务，就会从 app 走向 OS 位置。
\end{importantbox}

\subsection{中国御三家与模型 OS}

访谈提到中国御三家，但重点不在具体排名，而在一个结构判断：国内模型团队若要抓住第二幕，需要同时理解 Coding、Agent、成本和应用场景。模型 OS 的竞争不是只做一个聊天产品，而是接入办公、开发、内容、数据、自动化和企业流程。

\subsection{本章小结}

模型成为新一代 OS，意味着模型从回答系统变成任务系统。Coding 和 harness 是它进入这个位置的关键路径。

